% Вариант: Python-разработчик (backend).
% ATS-friendly resume: одна колонка, без таблиц, иконок и графики,
% стандартные заголовки разделов, контакты текстом, выделяемый текст.
\documentclass[10pt,a4paper]{article}

\usepackage[T2A]{fontenc}
\usepackage[utf8]{inputenc}
\usepackage[russian]{babel}
\usepackage[left=1.5cm,right=1.5cm,top=1.3cm,bottom=1.3cm]{geometry}
\usepackage{enumitem}
\usepackage[hidelinks]{hyperref}

\pagestyle{empty}
\hyphenpenalty=10000
\exhyphenpenalty=10000
\raggedright
\setlength{\parindent}{0pt}
\setlength{\parskip}{0pt}

\newcommand{\cvsection}[1]{%
  \vspace{7pt}{\large\bfseries #1}\par\vspace{2pt}\hrule\vspace{4pt}}

% #1 -- должность, #2 -- даты, #3 -- организация и город
\newcommand{\cventry}[3]{%
  \textbf{#1} \textbar{} #3 \textbar{} #2\par\vspace{2pt}}

\newlist{cvitems}{itemize}{1}
\setlist[cvitems]{label=\textbullet, leftmargin=1.2em, topsep=2pt, itemsep=1pt, parsep=0pt}

% --- Блоки навыков (порядок зависит от варианта) ---
\newcommand{\skilllangs}{\textbf{Языки программирования:} Python, SQL\par}
\newcommand{\skillai}{\textbf{Искусственный интеллект и ML:} большие языковые модели (LLM), LLM-агенты,
RAG (Retrieval-Augmented Generation), BM25, OpenAI API, LangChain, промпт-инжиниринг,
распознавание документов (OCR), A/B-тестирование (A/B testing), Claude Code\par}
\newcommand{\skillframeworks}{\textbf{Фреймворки:} FastAPI, SQLAlchemy (async), Alembic, Pydantic, pytest\par}
\newcommand{\skilllibs}{\textbf{Библиотеки:} asyncio, httpx, aiohttp, pandas, NumPy, scikit-learn, structlog\par}
\newcommand{\skilldb}{\textbf{Базы данных и очереди:} PostgreSQL, Redis, Apache Kafka\par}
\newcommand{\skillinfra}{\textbf{Инфраструктура и DevOps:} Docker, Kubernetes (K8s), Helm, GitLab CI/CD, Git, Linux\par}
\newcommand{\skillhuman}{\textbf{Языки:} русский --- родной, английский --- Intermediate (B1)}

\begin{document}

{\LARGE\bfseries Кирилл Козлов}\par\vspace{3pt}
Middle Python-разработчик (FastAPI, PostgreSQL, Kubernetes)\par\vspace{3pt}
Москва, Россия | Телефон: +7 (910) 968 1153 | Email: kirillandreevickozlov@yandex.ru\par
GitHub: github.com/noobmaster-ru | Telegram: @noobmaster\_rus

\cvsection{О себе}
%
Python-разработчик, бакалавр ВМК МГУ, учусь в магистратуре ВМК МГУ по машинному обучению.
В X5 Tech разработал три бэкенд-сервиса на FastAPI и развернул их в Kubernetes: LLM-агент для проверки
договоров (свыше 500 проверок на проде, эффект 8 FTE), обработка коммунальных платежей (эффект 10 FTE),
анализ иностранных инвойсов. Ключевые навыки: Python, FastAPI, PostgreSQL (SQLAlchemy, Alembic),
Apache Kafka, Docker и Kubernetes, GitLab CI/CD.%


\cvsection{Ключевые навыки}
\skilllangs
\skillframeworks\skilldb\skillinfra\skilllibs\skillai
\skillhuman

\cvsection{Опыт работы}
\cventry{Python-разработчик}{Апрель 2026 --- настоящее время}{X5 Tech, департамент ИИ-агентов, Москва}
\begin{cvitems}
%
    \item Разработал бэкенд-сервис на FastAPI и PostgreSQL (SQLAlchemy async, Alembic) для LLM-агента проверки договоров: свыше 500 проверок на проде за май--август 2026; экономический эффект для X5 Group --- 8 FTE.
    \item Сократил время подготовки юридического заключения до 8--19 минут против около 97 минут ручной работы юриста.
    \item Разработал бэкенд на FastAPI для сервиса обработки коммунальных платежей: приём результатов OCR из Apache Kafka, хранение и выдача во фронтенд; экономический эффект --- 10 FTE.
    \item Разработал и развернул на dev и prod stateless REST API для анализа иностранных инвойсов (OCR и LLM) для бизнес-единицы X5 Import; с августа 2026 в пилотной эксплуатации.
    \item Развернул сервисы в кластере Kubernetes (dev и prod): Helm-чарты, сборка и деплой образов в GitLab CI/CD, метрики латентности в Grafana.%

\end{cvitems}

\cvsection{Проекты}
\cventry{Python-разработчик, AxiomAI}{Октябрь 2025 --- Апрель 2026}{github.com/noobmaster-ru/axiomai}
\begin{cvitems}
    \item Разработал и развернул Telegram-бота на aiogram с использованием OpenAI API, PostgreSQL, Redis и Docker.
    \item Снизил нагрузку на менеджеров на 80\% и повысил скорость ответов в бизнес-коммуникациях.
\end{cvitems}
\vspace{4pt}

\cventry{Python-разработчик, WB Parser (Telegram Mini App)}{Июль 2025 --- Сентябрь 2025}{github.com/noobmaster-ru/tg-mini-app-backend}
\begin{cvitems}
    \item Разработал парсер поисковой выдачи Wildberries с фильтрацией по органической позиции товара.
    \item Упаковал продукт в Telegram Mini App с бэкендом и базой данных; привлёк 220 активных пользователей.
\end{cvitems}
\vspace{4pt}

\cventry{Python-разработчик, Placa}{Июль 2024 --- Декабрь 2024}{github.com/noobmaster-ru/wb\_tables}
\begin{cvitems}
    \item Разработал сервис аналитики в Google Sheets для продавцов Wildberries на Google Apps Script с интеграцией Wildberries API.
    \item Автоматизировал ежедневную загрузку данных: копирование шаблона таблицы, ввод API-токена, запуск с панели управления.
    \item Обеспечил около 300 продавцов ежедневными данными о рекламных расходах.
\end{cvitems}

\cvsection{Образование}
\cventry{МГУ имени М.В. Ломоносова, магистратура}{Сентябрь 2026 --- настоящее время}{Факультет ВМК, кафедра математических методов прогнозирования (ММП), машинное обучение, Москва}
\vspace{2pt}
\cventry{МГУ имени М.В. Ломоносова, бакалавриат}{Август 2022 --- Август 2026}{Факультет вычислительной математики и кибернетики (ВМК), Москва}
\begin{cvitems}
    \item Deep Learning School (DLS), направления Computer Vision и NLP --- Школа глубокого обучения МФТИ.
    \item Летняя школа по A/B-тестированию (A/B testing), X5 Tech --- лето 2026.
\end{cvitems}

\end{document}
